Now we introduce background on convex reformulations.
Convex reformulations re-write \cref{eq:non-convex-relu-mlp} as a convex
program by enumerating the activations a single neuron in the
hidden layer can take on for fixed \( Z \) as follows:
\[ \calD_Z = \cbr{D = \diag(\mathds{1}(Z u \geq 0)) : u \in \R^d}. \]
This set grows as \(|\calD_Z| \le O(r(n/r)^r)\), where \(r := \text{rank}(Z)\)
\citep{pilanci2020convexnn}.
Each ``activation pattern'' \( D_i \in \calD_{Z} \) is associated with
a convex cone,
\[
	\calK_i = \cbr{u \in \R^d : (2D_i - I) Z u \succeq 0}.
\]
If \( u \in \calK_i \), then \( u \) matches \( D_i \), meaning \( D_i Z u = \rbr{Z u}_+ \).
For any subset \( \tilde \calD \subseteq \calD_Z \),
the convex reformulation is,
\begin{equation}\label{eq:convex-relu-mlp}
	\begin{aligned}
		\!\!\! \min_{v, w} \, & L\Big(\!\!\sum_{D_i \in \tilde \calD}\!\! D_i Z (v_i \!-\! u_i), y\Big)\!+\!\lambda\!\sum_{D_i \in \tilde \calD}\!\norm{v_i}_2\!+\!\norm{u_i}_2 \\
		                      & \text{s.t.} \quad v_i, u_i \in \calK_i.
	\end{aligned}
\end{equation}
\citet{pilanci2020convexnn} prove that this program and \cref{eq:non-convex-relu-mlp}
are equivalent in the following sense:
if \( \tilde \calD = \calD_Z \) and \( m \geq m^* \) for some \( m^* \leq n + 1 \),
then the two programs have the same optimal value and every solution to the convex program
can be mapped to a solution of the non-convex training problem and vice-versa.
Given a solution \( \rbr{v^*, u^*} \), optimal weights
for the ReLU problem are given by
\begin{equation}\label{eq:convex-to-relu}
	\begin{aligned}
		W_{1i} & = v_i^*/ \sqrt{\norm{v_i^*}}, \quad w_{2i} = \sqrt{\norm{v_i^*}}
		\\
		W_{1j} & = u_i^*/ \sqrt{\norm{u_i^*}}, \quad w_{2j} = -\sqrt{\norm{u_i^*}},
	\end{aligned}
\end{equation}
where we use the convention that \( 0/0 = 0 \).

In practice, learning with \( \calD_{Z} \) is intractable except when the data
are low rank.
\citet{mishkin2022convex} provide refined conditions on \( \tilde \calD \)
which are sufficient for \cref{eq:convex-relu-mlp} to be equivalent to the
non-convex problem, while \citet{wang2021hidden} show that
the minimum of every sub-sampled convex program is a stationary point of the
ReLU training problem.

\subsection{Gated ReLU Networks}

An alternative is the gated ReLU activation function,
\[
	\phi_{g}(Z, u) = \text{diag}(\mathds{1}(Z g \geq 0)) Z u,
\]
where \( g \in \R^d \) is a ``gate'' vector, which is also optimized.
The gated ReLU activation modifies the ReLU activation to
decouple the thresholding operator from the neuron weights.
Two-layer gated ReLU networks predict as follows:
\begin{equation}\label{eq:non-convex-gated-relu-mlp}
	h_{W_1, w_2}(Z) = \sum_{i=1}^m \phi_{g_i}(Z, W_{1i}) w_{2i}.
\end{equation}
\citet{mishkin2022convex} show that this gated ReLU neural network has
the convex reformulation,
\begin{equation}\label{eq:convex-gated-relu-mlp}
	\min_{u} \; L\Big(\sum_{D_i \in \tilde \calD} D_i Z w_i, y\Big)
	+ \lambda \sum_{D_i \in \tilde \calD} \norm{w_i}_2,
\end{equation}
where decoupling the activations from the neuron weights
allows \( u_i, v_i \in \calK_i \) to be merged.
The solution mapping for \( w^* \) and conditions for
for the convex program to be equivalent to
\cref{eq:non-convex-gated-relu-mlp} are similar to the ReLU case.

%\begin{equation}\label{eq:convex-to-gated}
%    \begin{aligned}
%        W_{1i} & = v_i^*/\sqrt{\norm{v_i^*}}, \quad w_{2i} = \sqrt{\norm{v_i^*}},
%    \end{aligned}
%\end{equation}
%and conditions for the convex Gated ReLU problem to be equivalent to
%\cref{eq:non-convex-gated-relu-mlp} are similar to the ReLU case.


